\section{Related Work and Motivation}
\label{sec:related}

\subsection{Bone Suppression and CT-Derived Supervision}
\label{sec:rw-bone}
Early learning-based rib suppression used massive-training artificial neural networks \cite{suzuki2006mtann}. Later work used
convolutional networks, for example a cascade of multi-scale networks in the gradient domain \cite{yang2017cascade} and other deep models
\cite{gusarev2017deep}. Recent methods train diffusion models on paired radiographs and DES soft-tissue images
\cite{chen2024bsdiff,sun2026bsldm,sun2025gllcm,sun2026debonedit}. The SZCH-X-Rays data set was compiled for BS-LDM \cite{sun2026bsldm}
and made public with DeBoneDiT \cite{sun2026debonedit}; these methods train on hundreds of paired images, at up to $1024\times1024$
pixels. Paired DES images are scarce, however, and they come from a few scanners.

CT removes this bottleneck. Angaitkar et al. \cite{angaitkar2026anatomy} build digitally reconstructed radiographs from chest CT, project
bone, non-lung soft tissue and lung separately, and train models that predict one component. Their projections are released on the
Hugging Face Hub and are built from CT-RATE volumes \cite{hamamci2026ctrate}: the release holds \num{21887} projections whose volume names
are those of CT-RATE. We train on these projections and use the released weights as the main comparator. Our models differ in encoder and
resolution (ConvNeXt-V2-T at 512~px against SwinV2-B at 1024~px) and in sample budget, and they add the twin regulariser, the ensemble
and the audit. For the encoder we start from SwinV2-T \cite{liu2022swinv2} and screen ConvNeXt V2 \cite{woo2023convnextv2}, DINOv3
\cite{simeoni2025dinov3} and RAD-DINO \cite{perezgarcia2025raddino} under the same recipe: a modern convolutional encoder, a general
self-supervised vision transformer and a vision transformer pre-trained on chest radiographs.

Suppressed images also serve downstream models: they have been used to classify tuberculosis-consistent findings
\cite{rajaraman2021bone}, and readers and computer-aided detection have been compared with and without them
\cite{freedman2011lung,schalekamp2014cad}. A suppressed image can be given to a classifier alone, stacked with the radiograph as
channels, or as a separate input stream. The last two add computation, which a fair comparison has to match with a control of equal
capacity.

\subsection{Lesion Fidelity and Counterfactual Auditing}
\label{sec:rw-fidelity}
Learned image-to-image models can remove or invent content. Distribution matching in translation networks can drop or hallucinate lesions
\cite{cohen2018distribution}, deep reconstruction can be unstable to small structures \cite{antun2020instabilities}, and hallucinations of
learned reconstruction have been defined and measured \cite{bhadra2021hallucinations,tivnan2024hallucination}. Inserting known features
or signals into images is an established way to test learned MRI reconstruction \cite{cheng2020addressing} and the effect of learned
denoising on signal detection \cite{li2021assessing}. For bone suppression, simulated nodules were inserted to measure contrast after a
hand-built clavicle suppression method \cite{hogeweg2013suppression}, and a phantom with a simulated nodule at many positions was used to
measure contrast after bone suppression \cite{takagi2020dose}.

These studies measure lesion contrast. Contrast alone cannot separate a suppressor that keeps a lesion from one that suppresses weakly,
and a fidelity score computed over the whole image is dominated by normal anatomy. Our audit therefore measures the fraction of a known
additive lesion signal that a learned suppressor moves into its bone image, on real radiographs, at sites shared by every model compared,
and it reports this fraction next to suppression strength.

\subsection{Uncertainty Quantification and Selective Classification}
\label{sec:rw-uq}
Deep ensembles give a simple estimate of predictive uncertainty \cite{lakshminarayanan2017ensembles}, and split-conformal prediction turns
an error score into intervals with a coverage guarantee under exchangeability \cite{lei2018conformal,angelopoulos2023conformal}. We use
both for the bone image. The guarantee holds only when calibration and test images are exchangeable; a suppressor trained and calibrated
on CT-derived projections and applied to real radiographs breaks this assumption, and the size of the resulting loss of coverage has to be
measured rather than assumed.

For classification uncertainty we use conformal risk control of the false-negative rate \cite{angelopoulos2022crc}, class-conditional
(Mondrian) conformal sets \cite{vovk2005algorithmic} and selective classification with the area under the risk--coverage curve (AURC)
\cite{geifman2017selective,geifman2019bias}. Selective classification fits a screening setting: predictions with low uncertainty are
accepted, and the rest are referred to a reader.

\subsection{Cross-Modality Knowledge Transfer and Explainability}
\label{sec:rw-xm}
Knowledge distillation trains a student to match the tempered outputs of a teacher \cite{hinton2015distilling}, and cross-modal
distillation uses paired images to transfer supervision from a modality with rich signal to one without it \cite{gupta2016crossmodal}. For
the chest, CT-RATE pairs non-contrast CT volumes with findings extracted from their reports \cite{hamamci2026ctrate}, and the released
projections of its volumes \cite{angaitkar2026anatomy} supply a radiograph-like view of the same scans. This makes distillation from CT to
radiographs possible with public data alone; we found no study that tests whether it improves real-radiograph classification.

For multi-label chest radiograph classification we use a class-query decoder over spatial tokens, as in query-based multi-label heads
\cite{liu2021query2label,ridnik2023mldecoder}, so that the attention of each class on each input stream can be read out. Attribution maps
are computed with Grad-CAM++ \cite{chattopadhay2018gradcampp} and Integrated Gradients \cite{sundararajan2017axiomatic}. Gradient-based
heat maps such as Grad-CAM++ were developed for convolutional networks but are widely applied to transformers; their localisation can be
scored against expert annotations with the pointing game \cite{zhang2018topdown}. Hallucination has been defined and measured mainly for
learned image reconstruction \cite{bhadra2021hallucinations,tivnan2024hallucination}. For classifiers we test it as the appearance of a
finding under perturbations that add none, and as the response to counterfactual lesions.

Table~\ref{tab:literature} compares the studies closest to ours. Three gaps motivate the design: lesion erasure by learned suppressors has
not been measured against an exactly known signal; gains from suppressed inputs have not been compared against a capacity-matched
control; and the uncertainty, localisation and hallucination of CT-supervised radiograph models have not been audited together.

\begin{table*}[!t]
\centering
\caption{Comparison of studies on bone suppression, suppressed images as model inputs, lesion-fidelity evaluation and cross-modal
distillation with the present work. Check marks indicate whether a study reports each element: a counterfactual or simulated-lesion
audit (Lesion), a capacity-matched control (Cap.), uncertainty quantification or calibration (UQ), validation on external data (Ext.)
and explanation maps (XAI). \checkmark: reported; (\checkmark): qualitative only; --: not reported; n/a: not applicable to the study
design (no model trained by the authors, no input added to a model, or no clinical task). Rows marked $^{\dagger}$ were assessed from the
abstract because the full text is not openly accessible; all other rows from the full text.}
\label{tab:literature}
\footnotesize
\setlength{\tabcolsep}{3pt}
\begin{tabularx}{\textwidth}{@{}>{\raggedright\arraybackslash}p{2.6cm}>{\raggedright\arraybackslash}p{3.4cm}%
>{\raggedright\arraybackslash}p{3.3cm}>{\raggedright\arraybackslash}X*{5}{c}@{}}
\toprule
Study & Data & Method & Key result & Lesion & Cap. & UQ & Ext. & XAI\\
\midrule
Suzuki et al.$^{\dagger}$ \cite{suzuki2006mtann} & DES bone images for training; 118 validation and 136 independent test radiographs (14
 institutions) & Multiresolution massive-training ANN (MTANN) & Ribs and clavicles suppressed while nodule and vessel visibility was kept
 & -- & n/a & -- & \checkmark & n/a\\
Yang et al.$^{\dagger}$ \cite{yang2017cascade} & 504 two-exposure DES radiographs (404 training, 100 test) & Cascade of multi-scale CNNs
 in the gradient domain (CamsNet) & Soft-tissue PSNR 38.7\,dB, bone suppression ratio 83.8\,\%; other X-ray machines shown visually & --
 & n/a & -- & (\checkmark) & n/a\\
Chen et al. \cite{chen2024bsdiff} & Paired radiographs and DES soft tissue of 120 patients, one hospital (7:2:1 split) & Conditional
 diffusion model (BS-Diff) & Diffusion-based suppression with a clinical image-quality evaluation & -- & n/a & -- & -- & n/a\\
Sun et al. \cite{sun2026bsldm} & SZCH-X-Rays and JSRT (separate splits); Shenzhen set for a downstream test & Conditional latent
 diffusion (BS-LDM) & High-resolution suppression; clinical and downstream lesion-detection evaluation & -- & n/a & -- & \checkmark
 & n/a\\
Sun et al. \cite{sun2025gllcm} & SZCH-X-Rays (592/74/75) and JSRT (192/24/25) & Global-local latent consistency model (GL-LCM) &
 Suppression in 50 sampling steps, 8.54\,s against 84.62\,s for BS-LDM & -- & n/a & -- & -- & n/a\\
Sun et al. \cite{sun2026debonedit} & SZCH-X-Rays (released here) and JSRT; Asraf set for a downstream test & Depth-driven conditional
 bridge diffusion transformer (DeBoneDiT) & Public paired DES data; downstream classification on an external set & -- & n/a & -- &
 \checkmark & n/a\\
Angaitkar et al. \cite{angaitkar2026anatomy} & \num{21887} CT-RATE projections; TBX11K, Node21, VinDr-CXR & Anatomy-decomposed CT
 projections, SwinV2-B at 1024~px & Detection gains and higher real-nodule contrast; no lesion-insertion audit & -- & n/a & -- &
 \checkmark & --\\
Hogeweg et al.$^{\dagger}$ \cite{hogeweg2013suppression} & 20 simulated radiographs; 253 radiographs; 100 with CT nodules simulated
 near the clavicles & Model-based suppression of elongated structures & Simulated nodules kept more contrast than the clavicles after
 suppression (1.35 vs.\ 0.55) & \checkmark & n/a & -- & n/a & n/a\\
Takagi et al.$^{\dagger}$ \cite{takagi2020dose} & Anthropomorphic phantom, 12\,mm simulated nodule at 66 positions & Commercial
 bone-suppression processing at low tube voltage & Nodule contrast and CNR rose away from the ribs; lower dose possible & \checkmark &
 n/a & -- & n/a & n/a\\
Freedman et al.$^{\dagger}$ \cite{freedman2011lung} & Reader study, 368 patients, 15 radiologists & Rib and clavicle suppression software
 & Reader LROC AUC rose from 0.460 to 0.558 & -- & n/a & n/a & n/a & n/a\\
Schalekamp et al.$^{\dagger}$ \cite{schalekamp2014cad} & Reader study, 300 radiographs from four institutions, 8 readers & CAD with
 bone-suppressed images available & CAD raised reader AUC from 0.812 to 0.841 beyond suppressed images & -- & n/a & n/a & n/a & n/a\\
Rajaraman et al. \cite{rajaraman2021bone} & JSRT and NIH-CC DES (suppression); Shenzhen and Montgomery TB sets & Suppressed image
 replacing the radiograph as classifier input (VGG-16) & TB AUC 0.899 to 0.954 (Shenzhen) with suppressed inputs; cross-institutional
 suppression test & -- & n/a & -- & \checkmark & \checkmark\\
Gupta et al. \cite{gupta2016crossmodal} & NYUD2 (RGB and depth), JHMDB (RGB and optical flow) & Cross-modal distillation of mid-level
 features & Depth-only detection mAP 34.2 to 41.7\,\% on NYUD2 & n/a & n/a & -- & n/a & --\\
\midrule
This work & Qure.ai/CT-RATE projections; SZCH-X-Rays, VinDr-RibCXR, JSRT, Node21, TBX11K, VinDr-CXR, NIH ChestX-ray14, CT-RATE;
 BDCXR-3257 & Twin lesion-invariance regulariser, $K=3$ ensemble; multi-stream classifier; 2.5D Swin CT teacher distilled into Swin X-ray
 models & $A=0.051$ vs.\ 0.149 (released) and 0.178 (control); suppressed inputs and CT distillation did not raise NIH macro-AUROC
 & \checkmark & \checkmark & \checkmark & \checkmark & \checkmark\\
\bottomrule
\end{tabularx}
\end{table*}
